\paragraph{Effect of partition thickness.}
A distinctive feature of the conjugate problem is that the partition reduces
heat transport even when it is of equal conductivity ($k_r=1$): the solid band
carries no convection, so thickening it replaces convective transport with pure
conduction over a wider strip. Doubling the partition to $T_p=0.2$ lowers the
grid-converged $\mathrm{Nu}$ from $1.059$, $1.751$ and $2.794$ to $1.020$, $1.432$ and
$2.123$ at $\mathrm{Ra}=10^{4}$, $10^{5}$ and $10^{6}$, reductions of $3.7\%$, $18.2\%$
and $24.0\%$. The reduction grows sharply with Rayleigh number: at
$\mathrm{Ra}=10^{4}$ transport is conduction-dominated and largely
thickness-insensitive, while at $\mathrm{Ra}=10^{6}$ convection carries most of the heat
and the wider stagnant band blocks proportionally more of it. This monotone, Rayleigh-amplified
thickness sensitivity reproduces the trend reported by \citet{Khatamifar2021-si},
who likewise find that $\mathrm{Nu}$ decreases with partition thickness and that the
effect is most pronounced at small $k_r$, where the partition cannot compensate
by conducting more efficiently.
